\documentclass[../libro]{subfiles}

\begin{document}

\ifSubfilesClassLoaded{\mainmatter\chapter{行列式}\clearpage}{}

\section{杂例}

\begin{example}
    设 \(A\) 是 \(n\)~级阵.
    设 \(k\) 是数.
    证明, \(\det {(kA)} = k^n \det {(A)}\).
\end{example}

这里, 我展现一个用行列式定义的方法.

不过, 先回想定义:

% subfile fix: repeat restatable
% see:
% determinanto - determinantoj.tex
\ifSubfilesClassLoaded{%
    \begin{restatable}[行列式]{definition}{DefinitionDeterminants}
        设 \(A\) 是 \(n\)~级阵 (\(n \geq 1\)).
        定义 \(A\) 的\emph{行列式}
        \begin{align*}
            \det {(A)}
            =
            \begin{dcases}
                [A]_{1,1},
                 & n = 1;    \\
                \sum_{i = 1}^{n}
                {(-1)^{i+1} [A]_{i,1} \det {(A(i|1))}},
                 & n \geq 2.
            \end{dcases}
        \end{align*}
    \end{restatable}
}{\DefinitionDeterminants*}

\begin{proof}
    我们用数学归纳法证明此事.
    取数 \(k\).
    具体地, 设 \(P(n)\) 为命题
    \begin{quotation}
        对任何 \(n\)~级阵 \(A\),
        \begin{align*}
            \det {(kA)} = k^n \det {(A)}.
        \end{align*}
    \end{quotation}
    则, 我们的目标是,
    对任何正整数 \(n\), \(P(n)\) 是对的.

    \(P(1)\) 是对的.
    任取 \(1\)~级阵 \(A = [a]\).
    那么, \(kA = [ka]\).
    所以,
    \begin{align*}
        \det {(kA)} = ka = k^1 \det {(A)}.
    \end{align*}

    现在, 我们假定 \(P(m-1)\) 是对的.
    我们要证 \(P(m)\) 也是对的.
    任取 \(m\)~级阵 \(A\).
    按 \(kA\)~的列~\(1\) 展开, 有
    \begin{align*}
        \det {(kA)}
        = {} & \sum_{i = 1}^{m}
        {(-1)^{i+1} [kA]_{i,1} \det {((kA) (i|1))}}
        \\
        = {} & \sum_{i = 1}^{m}
        {(-1)^{i+1} k [A]_{i,1} \det {((kA) (i|1))}}.
    \end{align*}
    注意, 每个 \((kA)(i|1)\) 都是 \(m-1\)~级阵.
    并且, 不难验证, \((kA)(i|1) = k(A(i|1))\).
    从而, 根据假定,
    \begin{align*}
        \det {((kA)(i|1))} = k^{m-1} \det {(A(i|1))}.
    \end{align*}
    所以
    \begin{align*}
        \det {(kA)}
        = {} & \sum_{i = 1}^{m}
        {(-1)^{i+1} k [A]_{i,1} k^{m-1} \det {(A(i|1))}}
        \\
        = {} & k k^{m-1}
        \sum_{i = 1}^{m}
        {(-1)^{i+1} [A]_{i,1} \det {(A(i|1))}}
        \\
        = {} & k^m \det {(A)}.
    \end{align*}
    所以, \(P(m)\) 是对的.
    由数学归纳法, 待证命题成立.
\end{proof}

当然, 本题也有其他的解法.
% 若您适当地运用行列式的性质,
% 您或许能较简单地解此题.

\vspace{2ex}

在讲后一个例前,
我有必要提到阵的新记号.

\begin{definition}
    设 \(A\), \(B\), \(C\), \(D\)
    分别是 \(m \times s\), \(n \times s\),
    \(m \times t\), \(n \times t\)~阵.
    我们约定, 记号
    \begin{align*}
        \begin{bmatrix}
            A & C \\
            B & D \\
        \end{bmatrix}
    \end{align*}
    表示一个 \((m + n) \times (s + t)\)-阵 \(M\),
    且对任何不超过 \(m + n\)~的正整数 \(i\)
    与不超过 \(s + t\)~的正整数 \(j\),
    \begin{align*}
        [M]_{i,j}
        = \begin{cases}
              [A]_{i,j},
               & \text{\(i \leq m\), 且 \(j \leq s\)}; \\
              [B]_{i-m,j},
               & \text{\(i > m\), 且 \(j \leq s\)};    \\
              [C]_{i,j-s},
               & \text{\(i \leq m\), 且 \(j > s\)};    \\
              [D]_{i-m,j-s},
               & \text{\(i > m\), 且 \(j > s\)}.
          \end{cases}
    \end{align*}
\end{definition}

比如, 设
\begin{align*}
     &
    A = \begin{bmatrix}
            10 & 12 & 14 \\
            11 & 13 & 15 \\
        \end{bmatrix},
    \\
     &
    B = \begin{bmatrix}
            16 & 20 & 24 \\
            17 & 21 & 25 \\
            18 & 22 & 26 \\
            19 & 23 & 27 \\
        \end{bmatrix},
    \\
     &
    C = \begin{bmatrix}
            28 & 30 & 32 & 34 & 36 \\
            29 & 31 & 33 & 35 & 37 \\
        \end{bmatrix},
    \\
     &
    D = \begin{bmatrix}
            38 & 42 & 46 & 50 & 54 \\
            39 & 43 & 47 & 51 & 55 \\
            40 & 44 & 48 & 52 & 56 \\
            41 & 45 & 49 & 53 & 57 \\
        \end{bmatrix}.
\end{align*}
则
\begin{align*}
    \begin{bmatrix}
        A & C \\
        B & D \\
    \end{bmatrix}
    =
    \begin{bmatrix}
        10 & 12 & 14 & 28 & 30 & 32 & 34 & 36 \\
        11 & 13 & 15 & 29 & 31 & 33 & 35 & 37 \\
        16 & 20 & 24 & 38 & 42 & 46 & 50 & 54 \\
        17 & 21 & 25 & 39 & 43 & 47 & 51 & 55 \\
        18 & 22 & 26 & 40 & 44 & 48 & 52 & 56 \\
        19 & 23 & 27 & 41 & 45 & 49 & 53 & 57 \\
    \end{bmatrix}.
\end{align*}

设 \(A\), \(B\), \(C\), \(D\)
分别是 \(m \times s\), \(n \times s\),
\(m \times t\), \(n \times t\)~阵.
再设
\begin{align*}
    M = \begin{bmatrix}
            A & C \\
            B & D \\
        \end{bmatrix}.
\end{align*}
不难验证, 若 \(i \leq m\), 且 \(j \leq s\),
则
\begin{align*}
    M(i|j)
    =
    \begin{bmatrix}
        A(i|j) & C(i|) \\
        B(|j)  & D     \\
    \end{bmatrix},
\end{align*}
其中, \(B(|j)\) 表示去除 \(B\)~的列~\(j\) 后得到的阵
(不去除它的行),
且 \(C(i|)\) 表示去除 \(C\)~的行~\(i\) 后得到的阵
(不去除它的列).

\begin{example}
    设 \(A\), \(D\) 分别是 \(n\), \(t\)~级阵.
    设 \(C\) 是 \(n \times t\)~阵.
    证明,
    \begin{align*}
        \det {
            \begin{bmatrix}
                A & C \\
                0 & D \\
            \end{bmatrix}
        }
        = \det {(A)} \det {(D)},
    \end{align*}
    其中, 左下角的 \(0\) 指 \(t \times n\)~零阵.
\end{example}

% 若您学过 ``按多列 (行) 展开行列式'',
% 再利用行列式的多线性
% (具体地, 若一个方阵含一列 \(0\), 或含一行 \(0\),
% 则它的行列式为 \(0\)),
% 本题是简单的;
% 不过, 为降低学习难度,
% 我标记它为选学内容.
% 故, 我还是用定义解此题.

我还是用定义解此题.

\begin{proof}
    我们用数学归纳法证明此事.
    取 \(t\)~级阵 \(D\).
    具体地, 设 \(P(n)\) 为命题
    \begin{quotation}
        对任何 \(n\)~级阵 \(A\), 任何 \(n \times t\)~阵 \(C\),
        \begin{align*}
            \det {
                \begin{bmatrix}
                    A & C \\
                    0 & D \\
                \end{bmatrix}
            }
            = \det {(A)} \det {(D)},
        \end{align*}
        其中, 左下角的 \(0\) 指 \(t \times n\)~零阵.
    \end{quotation}
    则, 我们的目标是,
    对任何正整数 \(n\), \(P(n)\) 是对的.

    \(P(1)\) 是对的.
    设 \(A = [a]\).
    记
    \begin{align*}
        J =
        \begin{bmatrix}
            A & C \\
            0 & D \\
        \end{bmatrix},
    \end{align*}
    其中, 左下角的 \(0\) 指 \(t \times 1\)~零阵.
    则
    \begin{align*}
             & \det {(J)}
        \\
        = {} & \sum_{i = 1}^{1+t}
        {(-1)^{i+1} [J]_{i,1} \det {(J(i|1))}}
        \\
        = {} &
        (-1)^{1+1} [J]_{1,1} \det {(J(1|1))}
        +
        \sum_{i = 2}^{1+t}
        {(-1)^{i+1} [J]_{i,1} \det {(J(i|1))}}
        \\
        = {} &
        1\, [A]_{1,1} \det {(D)}
        +
        \sum_{i = 2}^{1+t}
        {(-1)^{i+1} 0 \det {(J(i|1))}}
        \\
        = {} &
        [A]_{1,1} \det {(D)}
        \\
        = {} & \det {(A)} \det {(D)}.
    \end{align*}

    现在, 我们假定 \(P(m-1)\) 是对的.
    我们要证 \(P(m)\) 也是对的.
    任取 \(m\)~级阵 \(A\).
    记
    \begin{align*}
        M = \begin{bmatrix}
                A & C \\
                0 & D \\
            \end{bmatrix},
    \end{align*}
    其中, 左下角的 \(0\) 指 \(t \times m\)~零阵.
    则
    \begin{align*}
             & \det {(M)}
        \\
        = {} & \sum_{i = 1}^{m+t}
        {(-1)^{i+1} [M]_{i,1} \det {(M(i|1))}}
        \\
        = {} & \sum_{i = 1}^{m}
        {(-1)^{i+1} [M]_{i,1} \det {(M(i|1))}}
        +
        \sum_{i = m+1}^{m+t}
        {(-1)^{i+1} [M]_{i,1} \det {(M(i|1))}}
        \\
        = {} & \sum_{i = 1}^{m}
        {(-1)^{i+1} [A]_{i,1} \det {(M(i|1))}}
        +
        \sum_{i = m+1}^{m+t}
        {(-1)^{i+1} 0 \det {(M(i|1))}}
        \\
        = {} & \sum_{i = 1}^{m}
        {(-1)^{i+1} [A]_{i,1} \det {(M(i|1))}}.
    \end{align*}
    注意, \(1 \leq i \leq m\) 时,
    \begin{align*}
        M(i|1)
        = \begin{bmatrix}
              A(i|1) & C(i|) \\
              0      & D     \\
          \end{bmatrix},
    \end{align*}
    其中, 左下角的 \(0\) 指 \(t \times (m-1)\)~零阵.
    根据假定,
    \begin{align*}
        \det {(M(i|1))} = \det {(A(i|1))} \det {(D)}.
    \end{align*}
    从而
    \begin{align*}
        \det {(M)}
        = {} & \sum_{i = 1}^{m}
        {(-1)^{i+1} [A]_{i,1} \det {(A(i|1))} \det {(D)}}
        \\
        = {} & \left(\sum_{i = 1}^{m}
        {(-1)^{i+1} [A]_{i,1} \det {(A(i|1))}}\right)
        \det {(D)}
        \\
        = {} & \det {(A)} \det {(D)}.
    \end{align*}

    所以, \(P(m)\) 是对的.
    由数学归纳法, 待证命题成立.
\end{proof}

\end{document}

This material contains two exercises about determinant.
They are easy, and can be solved
only with the definition.
